\section{Descrição das Atividades Realizadas}
\label{sec:metodologia}

A metodologia foi organizada em etapas de aquisição, preparação e análise dos dados, seguidas pela execução de experimentos de reconstrução virtual de temperaturas. A Figura~\ref{fig:fluxograma_metodologia}
apresenta uma visão geral desse procedimento, destacando o processamento comum das séries temporais e as diferentes configurações experimentais avaliadas. Os experimentos complementares investigam a influência da quantidade e da seleção de sensores de referência, bem como a contribuição da umidade relativa para as estimativas.

\begin{figure}[p]
\centering
\begin{tikzpicture}[
  font=\small,
  >={Latex[length=2mm]},
  node distance=0.55cm,
  etapa/.style={draw=black!65,rounded corners=2pt,fill=blue!6,
    align=center,text width=10.1cm,minimum height=1.00cm,inner sep=5pt},
  ensaio/.style={draw=black!65,rounded corners=2pt,fill=green!7,
    align=center,text width=4.55cm,minimum height=1.75cm,inner sep=5pt},
  saida/.style={draw=black!65,rounded corners=2pt,fill=orange!10,
    align=center,text width=10.1cm,minimum height=1.00cm,inner sep=5pt},
  auxiliar/.style={draw=black!55,dashed,rounded corners=2pt,fill=gray!7,
    align=center,text width=10.1cm,minimum height=0.95cm,inner sep=5pt},
  fluxo/.style={-{Latex[length=2mm]},thick,draw=black!70}
]
\node[etapa] (fonte) {\textbf{Aquisição dos dados}\\QuestDB: ambiente 1006, nove sensores de temperatura};
\node[etapa,below=of fonte] (preparo) {\textbf{Pré-processamento}\\Agregação em 5 minutos e sincronização temporal};
\node[etapa,below=of preparo] (matriz) {\textbf{Matriz experimental}\\12.892 janelas completas de temperatura};
\node[etapa,below=of matriz] (divisao) {\textbf{Divisão cronológica}\\9.024 janelas de treinamento (70\%) e 3.868 de teste (30\%)};
\draw[fluxo] (fonte) -- (preparo);
\draw[fluxo] (preparo) -- (matriz);
\draw[fluxo] (matriz) -- (divisao);

\node[ensaio,below=1.35cm of divisao.south west,anchor=north west,xshift=0.02cm] (exp1) {\textbf{Comparação dos métodos}\\Média aritmética\\Regressão linear e Random Forest\\9 sensores};
\node[ensaio,right=0.65cm of exp1] (exp2) {\textbf{Quantidade de entradas}\\Regressão linear\\Correlação; $k\in\{1,2,3,4,6,8\}$\\9 sensores};
\node[ensaio,below=0.55cm of exp1] (exp3) {\textbf{Seleção espacial}\\Correlação versus\\proximidade\\6 sensores com coordenadas};
\node[ensaio,right=0.65cm of exp3] (exp4) {\textbf{Inclusão de umidade}\\Regressão linear\\Temperaturas versus\\temperaturas + umidades};
\begin{scope}[on background layer]
  \node[draw=black!45,dashed,rounded corners=3pt,fill=gray!3,
    fit=(exp1)(exp2)(exp3)(exp4),inner xsep=0.25cm,
    inner ysep=0.55cm] (grupo) {};
\end{scope}
\node[fill=white,inner sep=2pt,font=\bfseries\small] at ([xshift=-3.05cm]grupo.north) {Experimentos de reconstrução};
\draw[fluxo] (divisao.south) -- (grupo.north);

\node[saida,below=0.70cm of grupo.south] (avaliacao) {\textbf{Avaliação das reconstruções}\\Comparação com os valores observados\\MAE, RMSE e $R^2$};
\node[saida,below=of avaliacao] (analise) {\textbf{Análise e interpretação}\\Comparação por método, sensor-alvo e configuração};
\draw[fluxo] (grupo.south) -- (avaliacao.north);
\draw[fluxo] (avaliacao) -- (analise);

\node[auxiliar,below=0.5cm of analise] (energia) {\textbf{Auditoria independente dos registros de energia}\\Caracterização da disponibilidade; sem entrada nos modelos};
\end{tikzpicture}
\vspace{10pt}
\caption{Visão geral da preparação dos dados, dos experimentos de reconstrução e da avaliação. A análise espacial considera seis sensores com coordenadas conhecidas, e o experimento com umidade utiliza amostras pareadas. A auditoria de energia é independente dos modelos de reconstrução.}
\label{fig:fluxograma_metodologia}
\fautor
\end{figure}

\subsection{Operacionalização do Enriquecimento}
\label{subsec:operacionalizacao_enriquecimento}

A base foi organizada em três níveis. Primeiro, as consultas e a grade
temporal estruturaram as observações para análise conjunta. Em seguida,
os identificadores foram associados ao cadastro e às coordenadas
confirmadas, permitindo comparar critérios de seleção e representar
os sensores. Por fim, a umidade foi alinhada às temperaturas e incluída
nas entradas da regressão linear. Essas etapas são rastreáveis nos
notebooks e nas tabelas exportadas.

% O aprendizado supervisionado utiliza a temperatura do alvo como resposta histórica e as medições dos demais dispositivos como entradas. O experimento com umidade compara oito temperaturas com oito temperaturas e oito umidades, em amostras pareadas. Já as coordenadas são utilizadas no ordenamento por proximidade e no mapa, sem integrar a matriz de variáveis explicativas dos modelos individuais. Dessa forma, o uso espacial e a inclusão de outra grandeza ambiental são avaliados em experimentos complementares, e não como uma única configuração.

O aprendizado supervisionado utiliza a temperatura do sensor-alvo como resposta histórica e as medições dos demais dispositivos como variáveis de entrada. Para avaliar a contribuição da umidade relativa, foram comparadas duas configurações de regressão linear: a primeira utiliza as temperaturas dos oito sensores de referência, enquanto a segunda incorpora também suas respectivas medições de umidade, totalizando 16 variáveis explicativas. Ambas foram avaliadas sobre os mesmos registros temporais para cada sensor-alvo, permitindo comparar seus resultados em condições equivalentes.

As coordenadas espaciais, por sua vez, foram utilizadas no ordenamento dos sensores por proximidade e na elaboração do mapa térmico, sem integrar diretamente as variáveis explicativas dos modelos de reconstrução. Dessa forma, a utilização das informações espaciais e a inclusão da umidade foram investigadas em experimentos complementares, e não como uma única configuração.


\subsection{Delimitação do conjunto experimental}
\label{subsec:conjunto_experimental}

O experimento principal utilizou registros da tabela
\texttt{sensor\_temperature} do QuestDB associados ao ambiente
\texttt{1006}. Foram selecionados os sensores 1, 48, 887, 891, 86782,
86789, 86790, 86791 e 86793. A consulta abrangeu o intervalo de
10 de junho de 2026, às 00h00, inclusive, a 26 de julho de 2026,
às 00h00, exclusive, em UTC. Na exploração anterior, os nove sensores
apresentaram registros em todos os 46 dias desse intervalo, enquanto
o sensor 15 apresentou registros em 11 dias. A menor continuidade
desse dispositivo sustenta sua ausência do conjunto experimental,
sem implicar que os sensores selecionados tenham cobertura integral
em todas as janelas temporais.

A coleta exploratória pela API Smartlab, que inclui sensores
cadastrados nos grupos Cluster e Reunião, foi mantida separada deste experimento.
As temperaturas utilizadas nos experimentos e no mapa térmico provêm
exclusivamente da extração do QuestDB, sem combinar fontes
com convenções temporais ainda não confirmadas.

\subsection{Agregação, sincronização e valores ausentes}
\label{subsec:preprocessamento}

A consulta agrupou as temperaturas de cada sensor em intervalos de
cinco minutos, por meio de \texttt{SAMPLE BY 5m}, retornando a média
e a quantidade de medições por intervalo. Foram obtidas 116.438 linhas agregadas, cada uma correspondente a um par sensor--janela. Portanto, cada linha representa um intervalo temporal agregado para determinado dispositivo, e não uma medição individual.
A extração foi salva localmente em CSV, acompanhada da consulta SQL,
da identificação da fonte e de um resumo criptográfico SHA-256.

Os timestamps foram interpretados em UTC e os sensores organizados
em colunas. A matriz foi reindexada sobre a grade completa do período,
com 13.248 janelas. Foram verificadas a presença dos nove sensores,
a unicidade de cada par sensor--janela, a pertinência dos timestamps
ao período e à grade, e a ausência de valores infinitos.

Para comparar os métodos nas mesmas observações, foram mantidas
somente as janelas com temperatura disponível para todos os nove
sensores. Esse procedimento resultou em 12.892 linhas completas
(97,31\% da grade), com exclusão de 356 janelas. Não houve imputação,
remoção automática de temperaturas extremas ou imposição de um número
mínimo de medições por janela além da existência de uma média válida.
A contagem das medições foi preservada para auditoria. Assim, médias
baseadas em quantidades diferentes de registros participam da avaliação,
o que deve ser considerado na interpretação dos resultados.

\subsection{Métodos de reconstrução}
\label{subsec:metodos_reconstrucao}

Cada um dos nove sensores foi utilizado, alternadamente, como alvo.
Na comparação principal, os outros oito sensores forneceram as
entradas. A temperatura do alvo foi utilizada como resposta no
ajuste supervisionado e como referência na avaliação, sendo sempre
excluída das variáveis explicativas.

A média aritmética foi adotada como referência sem treinamento,
conforme a Equação~\ref{eq:media_sensores}. A regressão linear, descrita
na Equação~\ref{eq:regressao_linear}, foi ajustada por mínimos quadrados,
com intercepto, utilizando a implementação da biblioteca Scikit-learn
\cite{pedregosa2011scikit}. Os coeficientes foram estimados apenas no
treinamento, e a configuração completa foi exportada com a execução.

O terceiro método foi \texttt{RandomForestRegressor}, com 100 árvores,
semente aleatória 42 e processamento paralelo. Foram utilizados erro
quadrático como critério, amostragem com reposição, todas as variáveis
como candidatas em cada divisão, profundidade máxima sem limite
explícito, mínimo de duas amostras para divisão e uma por folha.
Não foi realizada otimização de hiperparâmetros.

\subsection{Protocolo temporal e métricas}
\label{subsec:protocolo_avaliacao}

A matriz completa foi ordenada cronologicamente e dividida sem
embaralhamento: as primeiras 9.024 linhas foram destinadas ao
treinamento e as 3.868 restantes ao teste. O treinamento compreendeu
10/06/2026 às 01h40 até 11/07/2026 às 15h55; o teste compreendeu
11/07/2026 às 16h00 até 25/07/2026 às 09h55, em UTC. As proporções
70\%/30\% foram aplicadas à quantidade de linhas após a exclusão
das janelas incompletas, não à duração contínua do período.

Todos os métodos foram avaliados nos mesmos timestamps, para os
mesmos nove alvos. Trata-se de reconstrução contemporânea: as entradas
são medições dos demais sensores na mesma janela do alvo. Como cada
entrada é uma média de cinco minutos, a avaliação corresponde à
reconstrução da janela agregada, e não à previsão antecipada de
uma temperatura futura.

As métricas principais foram MAE e RMSE, definidos nas
Equações~\ref{eq:mae} e~\ref{eq:rmse}, ambos em graus Celsius.
Como medida complementar, calculou-se
\begin{equation}
    R^2=1-\frac{\sum_{j=1}^{n}(T_j-\widehat{T}_j)^2}
                     {\sum_{j=1}^{n}(T_j-\overline{T})^2},
    \label{eq:r2}
\end{equation}
em que $\overline{T}$ é a média das observações do alvo no teste.
As métricas foram calculadas individualmente para cada alvo.
Os resumos apresentam suas médias aritméticas entre os nove alvos,
com pesos iguais. Em particular, a média dos RMSE não corresponde
ao RMSE calculado reunindo todas as previsões. O desvio entre alvos, quando exportado, descreve a variabilidade de desempenho entre sensores-alvo, sem constituir intervalo de confiança.

\subsection{Quantidade e seleção de sensores de entrada}
\label{subsec:quantidade_sensores}

Um segundo experimento utilizou regressão linear com
$k\in\{1,2,3,4,6,8\}$ entradas. Para cada alvo, os demais sensores
foram ordenados pela correlação de Pearson em ordem decrescente,
calculada exclusivamente no treinamento. Os primeiros $k$ sensores
desse ordenamento forneceram as entradas. A divisão temporal, os
nove alvos e as amostras de teste permaneceram iguais aos da
comparação principal. Com oito entradas, as métricas foram conferidas
contra as da regressão linear do primeiro experimento.

O procedimento permite descrever o efeito de variar $k$ com uma
regra de seleção fixa. Não compara diretamente essa regra com
seleção aleatória ou proximidade espacial. O menor erro observado
no teste é uma descrição das configurações avaliadas, não uma
estimativa independente do desempenho de uma configuração escolhida
com base nesse mesmo teste. A busca combinatória com divisão
60\%/20\%/20\%, presente no notebook exploratório, não integra
os resultados consolidados aqui.

\subsection{Análise espacial complementar e mapa térmico}
\label{subsec:metodologia_espacial}

% O arquivo \texttt{sensorsWithCoordinates} registra as coordenadas relativas de 14 dispositivos, confirmadas pelo autor. Desses, seis pertencem ao conjunto experimental do QuestDB: 48, 887, 891, 86790, 86791 e 86793. A análise complementar utiliza esses seis sensores como alvos e candidatos a entrada, preservando as mesmas 12.892 janelas completas e a divisão de 9.024 linhas de treinamento e 3.868 de teste. Os sensores 1, 86782 e 86789 não participam dessa análise por não possuírem coordenadas confirmadas no cadastro.

Conforme apresentado na Seção~\ref{subsec:ambiente_experimental}, foram mapeadas coordenadas relativas para 14 dispositivos, representados na Figura~\ref{fig:lablasdpcschem}. Entretanto, apenas seis desses dispositivos pertencem ao conjunto principal de nove sensores utilizado nos experimentos com dados do QuestDB: 48, 887, 891, 86790, 86791 e 86793.

A análise espacial complementar considera exclusivamente esses seis sensores, utilizados alternadamente como alvos e candidatos a entrada. Foram preservadas as mesmas 12.892 janelas completas e a divisão temporal de 9.024 observações para treinamento e 3.868 para teste. Os sensores 1, 86782 e 86789 não participam dessa análise por não possuírem coordenadas confirmadas.

Para cada alvo, foram comparadas duas regras de seleção de entradas:
correlação de Pearson decrescente no treinamento e distância euclidiana
crescente nas coordenadas relativas. Foram utilizados de um a cinco
sensores de entrada, com regressão linear e as mesmas métricas da
comparação principal. Os empates nos ordenamentos foram resolvidos
pelo menor identificador. Os valores de distância são expressos nas
unidades relativas do cadastro, sem atribuição de precisão métrica.
Com cinco entradas, a equivalência das métricas das duas regras foi
verificada numericamente.

A planta fornecida pelo autor indica uma divisória entre o cluster
e a região dos demais sensores, com uma abertura de comunicação à
direita. Para representação gráfica, a divisória horizontal foi
aproximada por $y=5{,}1$, entre $x=0$ e $x=5{,}7$, e a parede vertical
por $x=6{,}2$, acima de $y=5{,}1$. Essas posições foram inferidas da
figura e não constituem um levantamento de paredes. Os sensores do
cluster não foram candidatos nos experimentos espaciais; assim, a
seleção não atravessa a divisória para escolher entradas desse espaço.

O mapa térmico apresenta a média de cada um dos seis sensores nas
mesmas 12.892 janelas do período completo. Uma interpolação pelo
inverso da distância (IDW), com potência dois, foi utilizada apenas
para visualizar uma distribuição espacial estimada dessas médias.
Em cada ponto da grade, empregaram-se os seis valores observados;
na coincidência com um sensor, atribuiu-se diretamente seu valor.
O preenchimento foi limitado ao envoltório convexo dos sensores,
com exclusão da região do cluster. As posições dos outros dispositivos
foram indicadas sem temperaturas e sem preenchimento no cluster.

A IDW não foi avaliada como método de reconstrução neste experimento.
O mapa não representa uma medição contínua do campo térmico, não
resolve as trocas de calor através da abertura e não descreve uma
condição instantânea. As médias do período completo são usadas apenas
nessa visualização, sem participar da seleção das entradas ou do ajuste
dos modelos de regressão.

\subsection{Umidade como informação complementar}
\label{subsec:metodologia_umidade}

Uma extensão avaliou se a umidade dos sensores de entrada acrescenta
informação à reconstrução de temperatura. Os registros da tabela
\texttt{sensor\_humidity} foram consultados para os mesmos nove
identificadores, ambiente e limites UTC do experimento principal.
A grandeza foi tratada como umidade relativa (UR), em porcentagem,
coerente com os dispositivos DHT22 identificados no banco e com a
documentação técnica da família AM2302 \cite{aosongam2302}.
Essa identificação não substitui a verificação metrológica dos dispositivos.

A consulta retornou média, mínimo, máximo e contagens por sensor e
janela de cinco minutos. A matriz foi reindexada na mesma grade de
13.248 janelas. O critério de validade exclui integralmente janelas
sem valores numéricos ou com mínimo abaixo de 0\% ou máximo acima
de 100\%; nenhuma das janelas retornadas dos nove sensores foi
excluída por esse critério. Lacunas foram preservadas, sem imputação.

Foram comparadas duas regressões lineares com intercepto: uma com
as oito temperaturas dos demais sensores e outra com essas mesmas
temperaturas e suas oito umidades, totalizando 16 variáveis. Nem a
temperatura nem a umidade do alvo integram as entradas. Para cada
alvo, ambas as variantes foram ajustadas nas mesmas linhas de
treinamento e avaliadas nos mesmos timestamps de teste. A seleção
de casos completos considera apenas as umidades de entrada,
sem exigir a disponibilidade da umidade do alvo.

Foi mantida a fronteira temporal de 11/07/2026 às 16h00 UTC, sem
refazer o corte proporcional após a seleção. Todos os alvos mantiveram
9.024 linhas de treinamento. Para o alvo 48, foram preservadas as
3.868 linhas de teste; para os demais, foram retidas 3.866, pois duas
janelas não possuíam todas as umidades de entrada. A referência com
somente temperaturas foi reajustada nesse mesmo subconjunto, evitando
atribuir ao enriquecimento uma diferença decorrente das amostras.
As médias das métricas foram calculadas com pesos iguais entre alvos.

\subsection{Auditoria exploratória do campo de energia}
\label{subsec:metodologia_energia}

A tabela \texttt{sensor\_energy} foi consultada para o mesmo ambiente
e período do experimento principal, preservando os identificadores
dos dispositivos, os modelos registrados e os limites temporais.
Os dados disponíveis foram agregados em intervalos de cinco minutos
para caracterização e comparação temporal com as séries de temperatura.

Entre os registros consultados, o dispositivo de identificador 23
estava associado ao modelo ACS712, destinado à medição de corrente
elétrica \cite{allegroacs712}. Embora esse componente permita obter
medidas de corrente em ampères (A), não foi possível confirmar a
unidade física dos valores armazenados no campo \texttt{energy}
a partir das informações disponíveis sobre o processamento das
medições. Por esse motivo, os registros foram mantidos em sua escala
original, sem conversões para corrente, potência ou consumo energético.

Dessa forma, essa etapa foi limitada à caracterização da
disponibilidade dos registros e à comparação temporal exploratória
com as temperaturas observadas. Não foram realizadas inferências
sobre os estados de funcionamento dos equipamentos ou estabelecidas
relações causais entre os registros elétricos e as condições
térmicas do ambiente.

% O cadastro associa o sensor 23 ao dispositivo denominado \textit{Ar Samsung Direita Vilalba}; o modelo registrado é ACS712.

% O fabricante descreve o ACS712 como sensor de corrente \cite{allegroacs712}. Entretanto, não foi confirmada a transformação aplicada pelo firmware até o campo \texttt{energy} no QuestDB. Portanto, os valores são apresentados na escala original, com unidade não confirmada. Não foram convertidos para potência ou consumo, integrados no tempo, nem utilizados para classificar arbitrariamente o estado de funcionamento do equipamento. A sobreposição temporal é descritiva e não constitui avaliação causal da climatização.

%O modelo ACS712 é destinado à medição de corrente elétrica \cite{allegroacs712}. Entretanto, a unidade dos valores armazenados no campo \texttt{energy} não pôde ser determinada a partir das informações disponíveis nesta análise. Por esse motivo, os registros foram mantidos em sua escala original, sem conversão para corrente, potência ou consumo energético.

\subsection{Organização e rastreabilidade}
\label{subsec:reprodutibilidade}

O notebook \texttt{02\_preprocessamento.ipynb} produz as matrizes,
as contagens e a auditoria da preparação; o
\texttt{03\_experimentos.ipynb} lê a matriz local e exporta métricas,
previsões, figuras e configuração da execução. A execução consolidada
foi registrada em 30 de outubro de 2026, com Python 3.12.3,
NumPy 2.5.4, Pandas 3.0.6 e Scikit-learn 1.9.1.
O notebook \texttt{04\_analise\_espacial.ipynb} executa a comparação
espacial, o inventário das extrações locais e a geração do mapa.
O arquivo de configuração preserva as versões, os parâmetros,
os limites dos conjuntos e o SHA-256 da matriz utilizada.
As extensões de umidade e energia foram executadas em 31 de outubro
de 2026 nos notebooks \texttt{05\_umidade.ipynb} e
\texttt{06\_energia\_exploratoria.ipynb}, com consultas, caches
e configurações próprios, preservando os experimentos anteriores.

Os códigos utilizados na preparação dos dados, na execução dos
experimentos e na geração das visualizações estão disponíveis
no repositório público do projeto no GitHub\footnote{
\url{https://github.com/GuScarenci/temperatureMonitoring}}.
O repositório reúne os materiais de implementação e documentação
associados ao desenvolvimento deste trabalho, permitindo consultar
os procedimentos adotados e apoiar sua reprodução.
